% -----------------------------------------------------------------------------
% Poster theme and reusable building blocks
% -----------------------------------------------------------------------------
% This file centralizes styling. Content files should mainly contain text,
% tables, figures, and calls to the macros defined here.

\setbeamercolor{background canvas}{bg=PosterBackground}
\setbeamercolor{normal text}{fg=PosterInk,bg=PosterBackground}
\setbeamerfont{normal text}{size=\Large}

\newenvironment{posterpage}{%
  \begin{minipage}[t][\textheight][t]{\textwidth}%
  \vspace*{0mm}%
}{%
  \end{minipage}%
}

% Box style for normal content blocks.
\tcbset{
  posterbox/.style={
    enhanced,
    colback=white,
    colframe=PosterPrimary!78!black,
    boxrule=1.8pt,
    arc=4mm,
    left=\PosterBoxPadding,
    right=\PosterBoxPadding,
    top=7mm,
    bottom=7mm,
    before skip=0pt,
    after skip=0pt,
    fonttitle=\LARGE\bfseries,
    coltitle=PosterPrimary,
    attach boxed title to top left={xshift=3mm,yshift=-3mm},
    boxed title style={
      colback=white,
      colframe=white,
      boxrule=0pt,
      left=1mm,
      right=2mm,
      top=0mm,
      bottom=0mm
    }
  },
  highlightbox/.style={
    enhanced,
    colback=PosterSoft,
    colframe=PosterPrimary!90!black,
    boxrule=2.2pt,
    arc=4mm,
    left=10mm,
    right=10mm,
    top=8mm,
    bottom=8mm,
    before skip=0pt,
    after skip=0pt
  },
  footerbar/.style={
    enhanced,
    colback=PosterFooterBackground,
    colframe=PosterFooterBackground,
    boxrule=0pt,
    arc=0mm,
    left=10mm,
    right=10mm,
    top=7mm,
    bottom=7mm,
    width=\linewidth,
    before skip=0pt,
    after skip=0pt
  }
}

\newtcolorbox{PosterBlock}[2][]{posterbox,title={#2},#1}
\newtcolorbox{PosterHighlight}[1][]{highlightbox,#1}

\newcommand{\PosterLead}[1]{%
  {\LARGE\bfseries\color{PosterInk}#1}\par
}

\newcommand{\PosterSmallNote}[1]{%
  {\large\color{PosterMuted}\raggedright #1\par}
}

\newcommand{\PosterMetric}[2]{%
  \begin{minipage}[t]{0.31\linewidth}
    \centering
    {\Huge\bfseries\color{PosterPrimary}#1}\par\vspace{1mm}
    {\large\color{PosterInk}#2}\par
  \end{minipage}%
}

\newcommand{\PosterLogoPlaceholder}[1]{%
  \fbox{%
    \begin{minipage}[c][18mm][c]{0.86\linewidth}
      \centering {\large\color{PosterMuted}Logo}\par
    \end{minipage}%
  }%
}

\newcommand{\PosterLogo}[2][height=28mm]{%
  \IfFileExists{#2}{\includegraphics[#1]{#2}}{\PosterLogoPlaceholder{#2}}%
}

\newcommand{\PosterFigurePlaceholder}[2][55mm]{%
  \begin{tikzpicture}
    \draw[rounded corners=3mm, line width=1.2pt, color=PosterPrimary!55]
      (0,0) rectangle (0.985\linewidth,#1);
    \node[align=center, text width=.82\linewidth, color=PosterMuted, font=\Large]
      at (.4925\linewidth,.5*#1) {#2\\[1mm]\large Datei in \texttt{assets/figures/} ablegen};
  \end{tikzpicture}%
}

\newcommand{\PosterGraphic}[4][width=\linewidth]{%
  \begin{figure}
    \centering
    \IfFileExists{#2}{\includegraphics[#1]{#2}}{\PosterFigurePlaceholder{#4}}%
    \caption{#3}
  \end{figure}%
}

\newcommand{\PosterQrPlaceholder}{%
  \begin{tikzpicture}[baseline=(current bounding box.center)]
    \draw[rounded corners=2mm, line width=1pt, color=PosterPrimary!70]
      (0,0) rectangle (30mm,30mm);
    \foreach \x/\y in {4/4,4/20,20/4,10/10,16/16,21/22,23/12,12/23}{
      \fill[PosterPrimary!80] (\x mm,\y mm) rectangle +(4mm,4mm);
    }
    \node[align=center, color=PosterMuted, font=\small] at (15mm,-4mm) {QR-Platzhalter};
  \end{tikzpicture}%
}

\newcommand{\PosterFooterQrPlaceholder}{%
  \begin{tikzpicture}[baseline=(current bounding box.center)]
    \fill[white] (0,0) rectangle (34mm,34mm);
    \draw[rounded corners=1mm, line width=0.8pt, color=PosterFooterText]
      (0,0) rectangle (34mm,34mm);
    \foreach \x/\y in {4/4,4/23,23/4,10/10,16/16,24/24,25/13,13/25}{
      \fill[PosterFooterBackground] (\x mm,\y mm) rectangle +(4mm,4mm);
    }
  \end{tikzpicture}%
}

\newcommand{\PosterQrGraphic}[1][height=34mm]{%
  \IfFileExists{\PosterQrCode}{\includegraphics[#1]{\PosterQrCode}}{\PosterFooterQrPlaceholder}%
}

\newcommand{\PosterFooterItem}[2]{%
  {\Large\bfseries\color{PosterFooterText}#1}\par\vspace{1.5mm}%
  {\large\color{PosterFooterMuted}\raggedright #2\par}%
}

% The footer is drawn as a page overlay so it stays anchored to the bottom
% margin independent of how much placeholder content is currently on the page.
\newcommand{\PosterFooterBar}{%
  \begin{tikzpicture}[remember picture,overlay]
    \node[anchor=south, inner sep=0pt] at ([yshift=\PosterOuterMargin]current page.south) {%
      \begin{minipage}{\textwidth}
        {\color{PosterAccent}\rule{\linewidth}{2.2pt}}\par\vspace{0pt}%
        \begin{tcolorbox}[footerbar]
          \begin{minipage}[c][\PosterFooterHeight][c]{0.14\linewidth}
            \centering
            \PosterQrGraphic
          \end{minipage}%
          \hfill
          \begin{minipage}[c][\PosterFooterHeight][c]{0.25\linewidth}
            \PosterFooterItem{\PosterFooterLeftTitle}{\PosterFooterLeftText}
          \end{minipage}%
          \hfill
          \begin{minipage}[c][\PosterFooterHeight][c]{0.27\linewidth}
            \PosterFooterItem{\PosterFooterMiddleTitle}{\PosterFooterMiddleText}
          \end{minipage}%
          \hfill
          \begin{minipage}[c][\PosterFooterHeight][c]{0.25\linewidth}
            \PosterFooterItem{\PosterFooterRightTitle}{\PosterFooterRightText}
          \end{minipage}
        \end{tcolorbox}%
      \end{minipage}%
    };
  \end{tikzpicture}%
}

\newcommand{\PosterTag}[1]{%
  \tikz[baseline=(X.base)]\node[fill=PosterAccent!16, draw=PosterAccent!80!black,
    rounded corners=2mm, inner xsep=3mm, inner ysep=1.4mm, font=\large\bfseries,
    text=PosterInk] (X) {#1};%
}

\newcommand{\PosterRule}{\par\vspace{3mm}{\color{PosterPrimary}\rule{\linewidth}{1.5pt}}\vspace{3mm}\par}
